\cotefile{64}
\scanpage{131}
Modules des courbes elliptiques : théorie transcendante

La donnée d'une courbe elliptique \inserted{$E$} sur $\mathbb C$ revient à celle \deleted{\unsure{à}} d'un \inserted{$\mathbb C$-}espace vectoriel de rang $1$, savoir \deleted{$\unsure{\mathfrak t_E|t_E}$}
\[
 V=\mathfrak t_E \qquad \text{(espace tangent à l'origine)} \tag*{(1)}
\]
et d'un ``réseau''
\[
 H\subset V \tag*{(2)}
\]
(sous-\deleted{groupe abélien discret}\inserted{$\mathbb Z$-module discret, libre} de rang $2$). La donnée du \deleted{groupe} $\mathbb Z$-module $\Pi$ permet de récupérer $V$ en tant qu'espace vectoriel réel \inserted{(de rang $2$)} par
\[
 V\simeq H\otimes_{\mathbb Z}\mathbb R\mathrel{\underset{\unsure{\circ}}{\overset{\mathrm{def}}{=}}}H_{\mathbb R} \tag*{(3)}
\]
et $E$ en tant que groupe de Lie réel, par
\[
 E\simeq V/H\simeq H_{\mathbb R}/H\,, \tag*{(4)}
\]
tandis que la structure complexe de $E$ est déterminée par la donnée d'une structure d'espace vectoriel complexe sur $V$ par passage aux quotients.

Étant donné un espace vectoriel réel $V$\inserted{\deleted{de rang \unsure{$n$}}}, la donnée d'une structure complexe sur $V$ (i.e. d'un \unsure{autom|endom} $u$ de $V$ de carré $-\mathrm{id}_V$) équivaut aussi à celle d'un \inserted{$\mathbb C$-}sous-espace vectoriel $L$ de \deleted{$\unsure{V_{\mathbb C}}$}
\[
 V_{\mathbb C}=V\otimes_{\mathbb R}\mathbb C \tag*{(5)}
\]
tel que la projection canonique
\[
 \begin{aligned}
 V_{\mathbb C}&\longrightarrow V=V_{\mathbb R}\\
 z&\longmapsto \tfrac12(z+\bar z)=\Re z
 \end{aligned} \tag*{(6)}
\]
\deleted{soit} \inserted{induise} un isomorphisme d'espaces vectoriels réels.
\[
 \text{\deleted{Pour un}}\quad\varphi\colon L\xrightarrow{\sim}V\quad(\mathbb R\text{-iso}) \tag*{(7)}
\]
% archive sheet number: 130
